\section{결론}
\label{sec:conclusion}

본 논문은 공급망 자료를 숨긴 채 회계 규칙을 검증하고,
근거에 문제가 생기면 관련 기록의 다음 사용을 막는 zkDPP의
구성과 보장 조건을 제시하였다.
거래마다 실제 입력과 출력의 미래 소비 식별자를 암호화해 남긴다.
거래 증명은 이 자료가 실제 계산과 맞는지도 검사한다.
승인된 감사자는 암호문을 열어 이력을 따라가고,
아직 소비되지 않은 관련 기록 전체를 차단한다.
회계 증표의 소유자 생성과 별도 이전 증표를 통해 미래 식별자를
생성 때 남기고, 같은 이력을 비공개 검증과 사후 집행에 사용한다.

보안 분석은 정상 요청이 기대한 결과를 내는 정확성,
틀린 거래·추적·완료 보고를 받아들이지 않는 건전성,
공개 자료에서 숨긴 내용을 구별하기 어렵게 하는 프라이버시로
구성하였다. 이 결과에는 명시한 암호·원장·권한 가정이 필요하다.
프라이버시는 무작위 오라클 모형에서 참여자의 비밀·개방값이
공개되지 않은 정상 이력을 읽는 수동 외부 관찰자에 한정한다.

차단에서는 대상 전체를 하나의 공통 시점에 확인해야 한다.
각 대상의 동결 성공 기록이 있어도 모두가 동시에 차단된 적은
없을 수 있다. 또 조사한 뒤 다른 거래가 대상을 먼저 소비할 수 있다.
zkDPP는 이런 대상을 목록에서 지우지 않고 미완료 사유로 남긴다.
전체 대상이 미소비·차단 상태임을 확인하고 그 상태가 유지되는 동안에만
그 대상을 통한 후속 사용을 거부한다.

기호 모델에서는 네 사례군의 334개 추적 결과를 독립 계산과 대조했다.
두 계산은 일치했지만 기록에서 시작하는 규칙이 명세와 다른 부분이
있으므로 명세 전체를 검증한 결과는 아니다.
누락 자료, 부분 동결과 소비 경합도 검사하였다.
질량이 0인 출력에 탄소 잔여가 몰리는 사례는 계산식의 합계가
맞아도 최종 제품군의 탄소 배분까지 보장되지는 않음을 보여 주었다.

시제품에서는 분기 이력의 추적·동결·결과 확인을 연결한 대표 실행과
아홉 거래의 비용을 제시했다. 감사 기능은 측정 조건에서
제약 수를 31.70--202.93\%, gas를 9.69--54.67\% 증가시켰다.
Transfer에서는 암호 계산의 스칼라 곱셈과 감사 자료 저장이
각각 추가 회로 비용과 추가 gas의 큰 부분을 차지했다.
기록 1,000개의 후속 추적은 9.760초,
대상 1,000개를 개별 거래로 동결하는 데에는 60.540초가 걸렸다.
두 시간은 서로 다른 작업을 따로 측정한 값이다.
공정 입출력 수를 바꾼 구간에서는 출력 증가가 입력 증가보다
gas에 더 큰 변화를 동반하였다.

제시한 보장은 승인된 회계 정책에 따른 디지털 기록과 그 후속 사용에
관한 것이다. 입력 증빙의 진실성, 이미 발급된 제품 주장의 취소와
물리적 회수는 외부 절차의 영역이다.
후속 과제는 현행 암호 구성의 가정 충족과 전체 통합 실행을 확인하고,
측정 판본의 대응 및 독립 반복 통계를 확보하는 것이다.
